\documentclass[aps,prb,onecolumn,nofootinbib,superscriptaddress]{revtex4-2}
\usepackage{amsmath,amssymb,amsthm,mathtools,bm}
\usepackage{booktabs}
\usepackage[colorlinks=true,linkcolor=blue,citecolor=blue,urlcolor=blue]{hyperref}
\usepackage{microtype}
\begin{document}
\title{A spectral bound on relaxation of the trajectory-averaged circuit}
\author{Asadullah Bhuiyan}
\date{September 29, 2026}
\begin{abstract}
We obtain the two-point relaxation rate directly from an ordered product of
single-particle projectors. No trajectory sampling or explicit time evolution
is required. This rate bounds the full channel's relaxation rate from above.
\end{abstract}
\maketitle

We ask how quickly the circuit forgets its initial two-point function after
discarding the measurement record. We use the manuscript convention
$C_{ab}=\operatorname{Tr}(\hat\rho\hat c_a^\dagger\hat c_b)$; the canonical
CPU engine stores the centered matrix $G=2C-\mathbf{1}$.
For a normalized overcomplete Wannier (OW) mode $j$, let $P_j$ be its
projector in the stored correlation ordering and $Q_j=\mathbf{1}-P_j$.
Perfect feedback fills a lower-band mode ($s_j=1$) or depletes an upper-band
mode ($s_j=0$). Gain or loss contributes
$s_jP_j-\{P_j,\overline C\}/2$ to the covariance increment, while number
dephasing contributes $-\{P_j,\overline C\}/2+P_j\overline C P_j$.
The elementary channel $\mathcal E_j=\mathrm{Id}+\mathcal L_j$ therefore gives
\begin{equation}
 \overline C'=\overline C-\{P_j,\overline C\}
       +P_j\overline C P_j+s_jP_j
       =Q_j\overline C Q_j+s_jP_j.
\end{equation}
This is an exact finite reset, not a time-step approximation. Closure of
this equation does not require the averaged state to be Gaussian.

One cycle is the composition $\mathcal E_{j_M}\cdots\mathcal E_{j_1}$,
with the rightmost operation acting first. At cycle boundaries,
\begin{equation}
 \overline C_{n+1}=A\overline C_nA^\dagger+B,
 \qquad A=Q_{j_M}\cdots Q_{j_1},\qquad
 \delta C_{n+1}=A\delta C_nA^\dagger,
\end{equation}
where $B$ is independent of the input and
$\delta C=\overline C-C_{\mathrm{ss}}$. Column vectorization yields
$T=A^*\otimes A$. If $a_p$ are the eigenvalues of $A$, the covariance
multipliers are $a_pa_q^*$. Hence, provided $0<\rho(A)<1$,
\begin{equation}
 \rho(A)=\max_p|a_p|,\qquad
 \Delta_C=-\log\rho(T)=-2\log\rho(A)
       =\min_{p,q}\bigl[-\log|a_pa_q^*|\bigr].
\end{equation}
The largest survival factor gives the smallest decay rate, measured per
cycle. This is not the occupation-spectrum gap around $1/2$, nor the
largest singular value of $A$. Non-normality can affect transients;
Jordan blocks add polynomial prefactors without changing the exponential rate.

The adjoint channel preserves
$\mathcal V_2=\operatorname{span}\{\hat{\mathbf{1}},\hat c_a^\dagger\hat c_b\}$,
and its restriction has diagonal blocks $1$ and $T^{\mathsf T}$.
Spectral inclusion then implies, for a mixing full channel in the same
physical symmetry sector,
\begin{equation}
 r_{\mathrm{full}}\geq\rho(T),\qquad
 \Delta_{\mathrm{full}}=-\log r_{\mathrm{full}}\leq\Delta_C.
\end{equation}
Mixing means a unique stationary state and no other unit-modulus eigenvalues
in that sector; $r_{\mathrm{full}}$ excludes the stationary eigenvalue. A closing upper
bound forces slow relaxation, but a finite bound does not establish a
many-body gap: higher-order observables may relax more slowly.

We evaluate $N_x=20$, $N_y=20,24,28,32,36,40,44,50$, and $\alpha_1=1,3$,
with $\alpha_2=30$, $n_{\mathrm{shell}}=1$, hard-wall support truncation at
the inclusive slab $x=5,\ldots,15$, and all slabs active. Both directions
are periodic, the twist is zero, and the trial orbitals are the $X$ pair.
The fixed raster-$y$ schedule advances $y$ at fixed $x$, with within-cell
order $A+,A-,B+,B-$. We construct $A$ by complex128 local rank-one updates
and compute its full non-Hermitian spectrum. There is no sampled ensemble,
chosen initial covariance, temporal average, or fitted decay law.
\input{results_summary.tex}
\end{document}
